\section{Introduction}

Cell-free massive multiple-input multiple-output (MIMO) has emerged as a transformative physical layer paradigm for 6G wireless networks. By deploying a large number of distributed access points (APs) connected to a central processing unit (CPU) via fronthaul/backhaul links, cell-free architectures can provide unprecedented macro-diversity and mitigate inter-cell interference, which is particularly well-suited for the Open Radio Access Network (O-RAN) framework \cite{pereira2025pixel, dai2025two}. However, the practical realization of this distributed architecture is fundamentally bottlenecked by the limited capacity of the fronthaul/backhaul links connecting the APs to the CPU \cite{xie2026meta, danaee2024relevance}. In realistic deployments, these links suffer from severe, dynamically fluctuating bandwidth constraints. Consequently, transmitting high-resolution, full-precision signals from all APs to the CPU is practically infeasible, necessitating highly efficient signal compression and quantization strategies to preserve the system's spectral efficiency and detection reliability \cite{danaee2025relevance}.

To address the capacity limitations of fronthaul links, various quantization and compression techniques have been investigated. Traditional uniform and non-uniform quantization methods, such as the Lloyd-Max algorithm, are typically designed to minimize the mean squared error of the reconstructed signal \cite{wang2020information}. However, these conventional approaches are fundamentally agnostic to the instantaneous channel state information (CSI) and the ultimate downstream task objectives, such as multi-user signal detection. This task-agnostic nature leads to severe performance degradation under low bit-rate constraints, as limited backhaul resources are often wasted on transmitting channel noise or deeply faded signal components \cite{shlezinger2021deep, zou2022goal}. Recently, artificial intelligence (AI) and deep learning have been extensively applied to MIMO detection \cite{he2020model, wei2020learned, sun2020learning} and resource allocation \cite{gao2020resource, syahran2026solving}. While some works have explored deep learning for signal quantization \cite{arvinte2020eq, feys2025learning} and information bottleneck-based feedback \cite{ma2022augmented, cao2022information}, existing AI-based approaches often decouple the quantization process at the APs from the global detection process at the CPU. Such decoupled designs fail to achieve a truly ``task-oriented'' optimal bit allocation \cite{mostaani2022task}, where the strictly limited backhaul capacity must be dynamically reserved for the most critical demodulation information based on the global network state. Furthermore, standard graph neural network (GNN) based detectors \cite{guo2026cell} assume uniform precision across all nodes, lacking the capability to robustly aggregate features corrupted by heterogeneous, asymmetric quantization noise.

To bridge these technical gaps, this paper proposes IB-CellFree-GNN, a novel task-oriented successive refinement framework rooted in the Variational Information Bottleneck principle. Unlike existing decoupled schemes, the proposed framework jointly optimizes the distributed quantization at the APs and the centralized multi-user detection at the CPU under dynamic backhaul constraints. The explicit novel contributions of this paper are summarized as follows:
\begin{itemize}
    \item \textbf{Differentiable Successive Refinement Policy Network:} We design a lightweight, local-CSI-driven policy network deployed at each AP to dynamically infer the optimal quantization bit-width (e.g., 0, 2, or 4 bits). By leveraging the Gumbel-Softmax trick, the discrete bit-selection process is rendered differentiable, enabling end-to-end Quantization-Aware Training (QAT). This mechanism effectively discards irrelevant or deeply faded signals, strictly reserving backhaul resources for the most critical information.
    \item \textbf{Bit-Aware Heterogeneous GNN for Robust Aggregation:} We develop a novel bit-aware heterogeneous GNN at the CPU to perform spatial aggregation and multi-user interference cancellation. The GNN adaptively adjusts its attention weights based on the inferred bit-depth of the incoming AP features, inherently compensating for asymmetric quantization noise and preventing low-precision signals from degrading the global detection accuracy.
    \item \textbf{Extreme Robustness and Scalable Computational Complexity:} Comprehensive numerical evaluations demonstrate that the proposed framework achieves near-optimal full-precision detection performance while consuming an average of only 2.5 bits per payload. Furthermore, the system exhibits extreme robustness, maintaining highly reliable communication even under a 50\% link failure rate. Finally, the AI-driven decision module exhibits an $\mathcal{O}(1)$ computational complexity scaling with respect to the number of AP antennas, confirming its exceptional scalability for practical 6G deployments.
\end{itemize}

The remainder of this paper is organized as follows. Section II rigorously defines the system model, including the cell-free MIMO architecture, signal models, and the mathematical formulation of the backhaul-constrained detection problem. Section III details the proposed IB-CellFree-GNN framework, elaborating on the successive refinement policy network, the Gumbel-Softmax-based QAT, and the bit-aware heterogeneous GNN. Section IV presents extensive numerical results and physical insights to validate the performance, robustness, and complexity of the proposed method. Finally, Section V concludes the paper and discusses potential directions for future research.